% !TeX root = ../main.tex
% !TeX encoding = UTF-8

\part{The AI Governance Failure in Detail}

\chapter{The Four Elements of a Digital Alibi}

Part Two established the same gap in six industries across seven decades. The gap was not technical. It was architectural: a governance design that captured decisions in their commercial or clinical or operational output without capturing the information picture that produced them. The industries that closed the gap, aviation, financial services, medicine, Lloyd's, did so by building an architecture that captured the four elements this book has named as the Digital Alibi: Decision Identity, The Information Picture, The Human Oversight Record, and The Accountability Chain.

Part Three is not a continuation of the historical argument. It is the operational architecture for boards who have accepted the argument and need to act on it. This chapter defines the four elements with the precision that implementation requires. Each element is examined at three levels: what it means as a governance standard, what it requires technically and operationally, and what it looks like in a specific sector context. The sectors used for the worked examples, credit decision-making, insurance underwriting, employment screening, and public sector procurement, were chosen because they collectively represent the range of organisational contexts in which AI-assisted decisions are currently generating the liabilities Part Two described.

The chapter closes with the two requirements that are frequently misunderstood: the simultaneity requirement, all four elements must be present at the same output moment, and the absence of any single element is a complete Digital Alibi gap, and the retrievability requirement, which is the dimension that most governance architectures address last and that most regulators will assess first.

\scenebreak

\section{Element 1 Decision Identity}

Decision Identity is the foundational element. Before any other aspect of a Digital Alibi can be established, the decision itself must be precisely identifiable. This is not a matter of assigning a reference number to a governance event. It is a matter of establishing, in a form that survives independent forensic examination, four pieces of information that must all be captured at the output moment.

\textbf{1.1 System identity} The specific AI system that produced the output contributing to the decision, named, versioned, and referenced in a form that allows independent verification. A system name alone is not sufficient. The version deployed at the output moment must be captured, because a system that has been updated since the output moment is not the system that produced it. Version control at the decision level, not the deployment level, is the requirement.

\textbf{1.2 Operational context} The configuration in which the system was operating at the output moment: the parameters, the operational scope, and any constraints or modifications applied to the standard deployment. A system running with modified parameters for a specific use case is a different operational entity from the baseline deployment, even if the version number is the same.

\textbf{1.3 Output identity} What the system produced, the specific score, classification, recommendation, or flag, in the exact form in which it was presented to the decision-maker. The output must be captured as produced, not as subsequently described or summarised.

\textbf{1.4 Decision attribution} The causal connection between the system's output and the decision taken, not asserted after the fact, but documented at the output moment as part of the governance record. This requires that the governance architecture distinguishes between decisions that used the AI output as the primary basis, decisions that used it as one of several inputs, and decisions in which the AI output was overridden.

\section{Decision Identity in practice: consumer credit}

A consumer credit lender uses an AI scoring model as the primary basis for lending decisions on applications below a set credit threshold. The model produces a score and a recommended decision, approve, decline, or refer for manual review, for each application.

Decision Identity for each application requires: the model version deployed at the moment the application was processed (not the current model version, which may have been updated); the specific parameter configuration active at that moment, including any recent adjustments to threshold settings; the exact score and recommended decision produced for this application, in the form presented to the underwriter or automated decision engine; and the classification of how that output was used, whether the automated decision engine acted on the recommendation directly, or whether a human reviewer was involved and acted consistently with, inconsistently with, or independently of the model's recommendation.

Most consumer credit operations can produce parts of this record from system logs. Few can produce all four components, for a specific application, on demand, in a form that ties the specific score to the specific model version and configuration that produced it. The gap is typically in version control at the decision level: the organisation knows what model version is currently deployed, but cannot demonstrate what version was deployed on a specific date three months ago for a specific application.

\scenebreak

\section{Element 2 The Information Picture}

The Information Picture is the most operationally complex element and the one most consistently absent from existing governance architectures. It has two dimensions of equal legal significance, and most current governance frameworks address only the first.

The positive record is what the AI system was given: the data inputs, the model state, and the parameters that collectively defined the information environment in which the system produced its output. This is the dimension most organisations attempt to capture through model logging and data lineage tracking, with varying degrees of completeness.

The negative space is what the AI system was not given: the data it did not receive, the context it was not designed to consider, and the inputs it was structurally blind to. The negative space is as evidentially significant as the positive record. A credit scoring model that was not given rental payment history produced a creditworthiness assessment in the absence of information that would have been relevant. An insurance pricing model that was not given mental health condition data produced a risk assessment without that context. An employment screening model that was not given information about career gaps produced a candidate assessment with a structural blind spot. In each case, the negative space is where discrimination and bias claims concentrate, and it is the dimension that the information picture must capture contemporaneously.

\textbf{2.1 Positive record} The complete set of data inputs presented to the system for this specific output, including the data sources, the data vintage (the date to which the data was current), and any transformations applied between raw data and model input. Data vintage matters because a model given data current to a date three months before the output moment is operating on a different information picture from a model given current data.

\textbf{2.2 Model state} The training data vintage, the last validation date, and the performance characteristics of the model at the version deployed, not the current version. A model that has been retrained since the output moment has different performance characteristics; the information picture must reflect the model as it was, not as it subsequently became.

\textbf{2.3 Negative space} A contemporaneous record of the data the system was not given: what information sources were not connected, what data fields were absent or null, and what contextual factors were structurally outside the model's information environment. This record must exist at the output moment, not be reconstructed from data architecture documentation after the fact.

\textbf{2.4 Confidence and uncertainty} Where the system produces a confidence indicator, probability score, or uncertainty measure alongside its primary output, this must be captured as part of the information picture. A recommendation produced with low confidence is a different piece of information from the same recommendation produced with high confidence, and the governance record must reflect the distinction.

\section{The Information Picture in practice: Lloyd's specialty insurance underwriting}

A Lloyd's managing agent deploys an AI exposure assessment model to assist underwriters in pricing complex specialty risks. The model ingests structured exposure data from the submission, applies catastrophe model outputs from a licensed third-party provider, and produces a technical price recommendation with a confidence interval.

The Information Picture for each bound risk requires: the complete exposure data submitted by the broker, including any fields that were absent or incomplete; the catastrophe model version and the peril set it covered, and, critically, the perils it did not cover, which represent the negative space for that assessment; the technical price recommendation and confidence interval as produced; and a record of any data quality flags the system raised during processing, which represent the system's own documentation of its negative space at the output moment.

The negative space is where Lloyd's regulatory exposure concentrates. An underwriting decision made on a catastrophe model that does not cover a specific emerging peril, cyber-physical interaction, for instance, or climate transition risk, has a negative space that is material to the risk assessment. If that negative space is not captured at the output moment, the managing agent cannot subsequently demonstrate, to Lloyd's Corporation or to a reinsurer disputing a claim, what the model could and could not see at the moment the risk was priced.

\scenebreak

\section{Element 3 The Human Oversight Record}

The Human Oversight Record is the element that most governance architectures treat as the easiest to produce and that, in practice, is the most frequently deficient. The record that a human reviewed an AI output is not a Human Oversight Record. The record of what that human did with the review is the Human Oversight Record.

The distinction matters because the legal standard for human oversight is not presence but substance. The EU AI Act's Article 14 requires that human oversight enables those responsible for the system to detect and correct failures. The FCA's SM\&CR requires that Senior Managers take reasonable steps to discharge their oversight obligations. In both cases, the question is not whether a human was in the loop. It is whether the human in the loop had the capacity and the authority to exercise genuine oversight, and whether they did.

Performative oversight, a governance process in which a human reviewer is present but structurally unable to exercise genuine review, is not an oversight record. It is a liability record. A reviewing officer who processes three hundred AI-assisted decisions per hour, who has no independent access to the data the model used, and who has never in practice overridden a model recommendation, has not exercised oversight. They have provided a signature. The governance record of their signature is evidence of their presence, not of their oversight.

\begin{quote}
\emph{The record that a human reviewed an AI output is not a Human Oversight Record. The record of what the human did with the information, what they had the capacity to review, what authority they had to override, and what they decided, is the Human Oversight Record. The distinction is the difference between presence and oversight.}
\end{quote}

\textbf{3.1 Reviewer identity and authority} The named individual who reviewed the AI output, their role and qualifications, and the specific authority they held to override the system's recommendation at the time of review. Authority must be documented as current at the review moment, a reviewer whose override authority was granted six months ago but was subsequently modified or removed cannot be shown to have had that authority at the specific output moment.

\textbf{3.2 Information presented} What the reviewer was actually shown: the AI output in the form it was presented, the information the reviewer had independent access to, and whether the reviewer had access to the underlying data the model used or only to the model's output and recommendation. A reviewer who saw only the score and the recommendation was in a different information position from a reviewer who had access to the input data and the model's confidence interval.

\textbf{3.3 Review duration and capacity} The time available for the review and the volume of decisions being reviewed at the same period. A reviewer with genuine capacity to examine the AI output for this specific decision is in a different oversight position from a reviewer processing decisions at a rate that structurally prevents meaningful examination of any individual decision. Capacity is part of the oversight record because it is part of the evidence of whether genuine oversight was possible.

\textbf{3.4 The override record} Whether the reviewer acted consistently with, inconsistently with, or independently of the AI system's recommendation, and the basis on which they did so. An override without a recorded basis is an override without an oversight record. The governance record must capture what the reviewer decided and why, at the decision moment.

\section{The Human Oversight Record in practice: employment screening}

A financial services firm uses an AI candidate screening system to rank applicants for graduate roles. The system processes applications and produces ranked shortlists with individual candidate scores. A human recruiter reviews the shortlist and confirms or modifies it before candidates are invited to interview.

The Human Oversight Record for each hiring cohort requires: the recruiter's name, their role in the hiring process, and their authority to modify the shortlist independently of line manager approval; what the recruiter was shown, whether they saw candidate scores alongside CVs, or only the ranked list without underlying scores; the time available for their review of the full shortlist; and for any modification to the shortlist, the documented basis for the modification.

The deficiency most likely to appear in an employment tribunal challenge to this process is in the override record. If the recruiter reviewed the shortlist and confirmed it without modification, the governance record typically shows ``shortlist confirmed.'' It does not show what the recruiter considered, what information they had access to, or whether they had the practical capacity to examine individual candidate assessments at the level the equality legislation requires. The Human Oversight Record must be specific enough to demonstrate that the oversight was genuine, that the recruiter was in a position to identify a candidate the model had undervalued and had both the information and the authority to add them.

\scenebreak

\section{Element 4 The Accountability Chain}

The Accountability Chain is named individual accountability from the AI system to the board, documented at the decision moment, not assigned after. It is the element most frequently confused with an organisational chart and the one whose confusion with an organisational chart most reliably produces a Digital Alibi gap.

An organisational chart describes who is responsible for what in structural terms. It may identify a Chief Risk Officer as accountable for model risk, a CTO as accountable for technology systems, and a board Risk Committee as accountable for risk appetite. This structural accountability is real and necessary. It is not an Accountability Chain at the decision level. The Accountability Chain requires that for this specific AI-assisted decision, on this specific date, the named individual accountable for each governance function, model ownership, decision process ownership, risk acceptance, is documented in the governance record as current at the output moment.

The distinction matters because the SM\&CR operates on named individual accountability at the decision moment, not on structural accountability at the organisational level. A managing agent whose SM\&CR accountability map names a Chief Underwriting Officer as accountable for AI-assisted underwriting decisions has, if that individual resigned three weeks before a specific decision was made and their replacement was not yet formally registered on the SM\&CR form, a gap in the Accountability Chain for that decision. The structural accountability remained in place. The named individual accountability at the decision moment did not.

\textbf{4.1 Model owner} The named individual accountable for the AI system's design, validation, and performance, specifically accountable for the version deployed at the output moment. This must be a current named individual, not a team or a function. Where model ownership transfers between individuals, through departure, role change, or interim appointment, the transfer must be documented, and the individual accountable at each specific output moment must be identifiable from the contemporaneous record.

\textbf{4.2 Decision process owner} The named individual accountable for the governance process through which the AI system's output is incorporated into the decision, the process that defines how the output is presented, what human review it receives, and what authority governs the final decision. This is distinct from the model owner: the decision process owner governs how the output is used; the model owner governs what the output is.

\textbf{4.3 Risk acceptance authority} The named individual at board level who has accepted the risk that this AI system represents, documented at the level of the specific risk, not at the level of the technology estate generally. This is the connection between the decision-level accountability and the board-level accountability that SM\&CR requires, and it must be documented at the decision moment for the chain to be complete.

\textbf{4.4 Chain currency} All three named individuals must be current at the output moment, not current at the date of the last annual review of the accountability map. A chain that is accurate for eleven months of the year and outdated for one, because of a departure, an interim appointment, or a role restructure, has a gap in the one month that a specific challenged decision may have occurred. Currency means documented contemporaneously, not confirmed retrospectively.

\section{The Vanishing Vendor: When the Chain Has No Internal Anchor}

The Accountability Chain has a specific vulnerability that becomes acute when AI systems are procured from third-party vendors rather than built and operated in-house. Chapter 5 established the independence design principle: the Digital Alibi architecture must be independent of the system it documents, and must survive vendor insolvency, data migration, and platform decommissioning. The Accountability Chain element exposes why independence is not merely a survivability requirement. It is an accountability requirement.

A government agency deploys an AI system for welfare benefit eligibility assessment, procured from a specialist software provider. The system ingests applicant data, applies a scoring model, and produces eligibility recommendations that an assessment officer reviews before a final decision is recorded. The Accountability Chain, as the agency has documented it, names the software provider as accountable for model ownership, the agency's digital team as accountable for the decision process, and a senior operational manager as the risk acceptance authority at board level.

Two years into deployment, the software provider enters administration. Its technical staff disperse. Its servers are transferred to an insolvency practitioner. The scoring model's design documentation, training data, and version history exist on the provider's infrastructure, which the agency cannot access. An applicant challenges an eligibility decision made fourteen months ago, citing potential bias in the assessment model. The agency must produce the Accountability Chain for that decision: the named model owner, the model version deployed at the output moment, and the basis on which the model owner accepted accountability for the version's performance characteristics.

The agency cannot produce it. The model owner, the software provider, no longer exists in any form that can be held accountable. The model version documentation is on infrastructure the agency does not control. The accountability chain named a vendor as the anchor of its model ownership. When the vendor vanished, the chain had nothing to attach to. The assessment officer's review exists in the agency's case management system. The three remaining elements of the Digital Alibi do not.

\begin{quote}
\emph{An accountability chain that names a vendor as its anchor is a chain that snaps when the vendor goes. The independence requirement, that the Digital Alibi architecture is owned, operated, and survivable independently of the system it documents, is not a technical nicety. It is the condition on which the chain remains intact when it is most needed.}
\end{quote}

The practical requirement this scenario establishes is that the Accountability Chain must have an internal anchor at every link. Where a third-party vendor provides the AI system, the organisation deploying it must name an internal individual as accountable for model oversight, not for building the model, but for verifying and accepting the vendor's model at the version deployed at each specific output moment. The vendor's documentation supports that acceptance. The named internal individual's contemporaneous acceptance record is the accountability chain link that must survive the vendor's departure.

\section{The Accountability Chain in practice: public sector procurement}

A local authority uses an AI scoring system to assess supplier bids for a framework contract. The system scores bids against weighted criteria and produces a ranked recommendation. A procurement officer reviews the ranked list and confirms the selection, subject to approval by a senior commissioning manager.

The Accountability Chain for each procurement decision requires: the named individual responsible for the AI scoring system and its current configuration, which may be a technology team member or a procurement systems manager; the named procurement officer who owns the governance process through which the AI recommendation is reviewed and confirmed; and the named senior commissioning manager who accepted the risk of AI-assisted procurement at the level of this framework contract.

The gap most likely to appear in a procurement challenge is in the currency of the chain during a change in commissioning management. If the senior commissioning manager who approved the use of the AI scoring system for this framework contract left the local authority after that approval and before the specific bid assessment being challenged, the authority must be able to demonstrate that their replacement had accepted the same risk acceptance for this framework at the moment the challenged decision was made, not that the governance structure provides for a commissioning manager to be accountable generally, but that a named individual had accepted the specific risk at the specific moment.

\scenebreak

\section{The Simultaneity Requirement: All Four, at the Same Moment}

The four Digital Alibi elements are not a checklist. They are not a scoring system in which three out of four produces a partial Digital Alibi. They are a simultaneous evidentiary standard: all four must be present at the output moment, in a form that is contemporaneous and retrievable, or the Digital Alibi does not exist.

The simultaneity requirement is frequently misunderstood in governance implementation. Organisations that have made progress on model versioning and data lineage, addressing Decision Identity and the positive record of The Information Picture, will sometimes describe themselves as ``on the way'' to a Digital Alibi. They are on the way to having two elements. They have not yet built any part of The Human Oversight Record or the Accountability Chain at the decision level. An organisation with two elements and two gaps has not achieved fifty per cent of a Digital Alibi. It has achieved nothing that will answer the output-moment question that a regulatory examination will ask, because that question requires all four elements simultaneously.

\begin{quote}
\emph{An organisation with three elements and one gap has not achieved seventy-five per cent of a Digital Alibi. It has achieved a governance record that will fail the output-moment question at the point of the missing element. The question does not give partial credit.}
\end{quote}

The practical implication of the simultaneity requirement is that Digital Alibi implementation must be designed as an integrated system, not as a set of independent governance improvements. Model versioning improvements that are not connected to the human oversight record produce a better information picture without the oversight evidence. SM\&CR accountability map updates that are not connected to the decision-level accountability chain produce a better structural record without the decision-moment attribution. The four elements must be captured in the same governance event, the output moment, or the simultaneity requirement is not met.

\scenebreak

\section{The Retrievability Requirement: What ``Preserved'' Actually Means}

The retrievability requirement is the dimension of the Digital Alibi standard that most governance architectures leave to last and that most regulators and litigants will assess first. An information picture that exists in archived logs that cannot be accessed within the timeframe a regulatory examination requires is not a Digital Alibi. It is an archive.

The retrievability requirement has three components that must all be met. First, the four elements must be stored in a form that can be accessed and presented to an examiner without manual reconstruction. A Digital Alibi that requires three weeks of data engineering to assemble from disparate system logs is not a Digital Alibi that can be produced on demand. Second, the integrity of the stored record must be verifiable, it must be demonstrable that the record has not been modified since the output moment at which it was created. Third, the record must be accessible within the timeframe that the specific regulatory or legal context requires, which varies by sector but which will, in most cases, be measured in days rather than weeks.

\begin{longtable}[]{@{}
  >{\raggedright\arraybackslash}p{(\columnwidth - 0\tabcolsep) * \real{1.0000}}@{}}
\toprule\noalign{}
\begin{minipage}[b]{\linewidth}\raggedright
\textbf{THE DIGITAL ALIBI RETRIEVABILITY TEST, FOUR QUESTIONS FOR ANY BOARD}
\end{minipage} \\
\midrule\noalign{}
\endhead
\bottomrule\noalign{}
\endlastfoot
1. Can we produce the complete Digital Alibi, all four elements, for any specific AI-assisted decision in our material decision processes, on demand, within five working days? If the honest answer is no, the retrievability requirement is not met. \\
2. Can we demonstrate that the record we produce has not been modified since the output moment? If the record is stored in a system that does not maintain an immutable audit trail, the integrity requirement is not met. \\
3. Does the record we produce reflect what the system knew at the output moment, or what the system currently knows? If the record would be assembled from current system states rather than from contemporaneous capture, it is a reconstruction, not a Digital Alibi. \\
4. Is the retrieval process documented and testable by a person who was not involved in building it? If the retrieval depends on institutional knowledge held by specific individuals, the survivability requirement, established in Chapter 5, is not met. \\
\end{longtable}

The retrievability test is the board-level governance question for every AI system in a material decision process. It does not require the board to understand the technical architecture of the retrieval system. It requires the board to receive confirmation, from a source that is independent of the teams responsible for the system, that the test can be passed. A board that has not received that confirmation does not know whether it has a Digital Alibi or an archive.

\scenebreak

\section{Identifying the Decisions to Which This Standard Applies}

Chapter 11 has defined the four elements and their operational requirements. The standard is now clear. The question that follows is which of an organisation's AI-assisted decisions are subject to it.

The answer is not ``all AI-assisted decisions.'' The Digital Alibi standard applies to material AI-assisted decisions, those whose outputs contribute to decisions with consequential implications for individuals, organisations, or regulatory obligations. Most organisations deploying AI are making both material and non-material decisions with AI assistance, and the governance burden is different for each category. Chapter 12 establishes a precise definition of materiality in the AI decision context, provides a decision classification framework that organisations can apply to their current AI deployments, and examines the specific challenge of AI decisions that appear routine but whose materiality becomes apparent only when their aggregate effect is examined.
